%% kapa: exam=2 round=31 year=2020 period=3 subject=compensation_law subject_ko=감정평가 및 보상법규 seq=2 number=2 points=30
%% source: Q-Net 감정평가사 기출문제 (cst00309) artlSeq=5207673 file=raw/1008/5207673/3교시_감정평가 및 보상법규.pdf
%% conversion: kordoc 4.24.51 -> tools/gen2.py (auto); review=proofread-gpt-astra; flags=none
\begin{longquestion}{【문제 2】}
A시의 시장 甲은 2018. 5. 31. 乙ㆍ丙 공동소유의 토지 5,729 m\textsuperscript{2}(이하 ‘이 사건 토지’라고 한다)에 대하여 2018. 1. 1. 기준 개별공시지가를 m\textsuperscript{2}당 2,780,000원으로 결정ㆍ고시하였다. 乙은 2018. 6. 19. 甲에게 「부동산 가격공시에 관한 법률」 제11조에 따라 이 사건 토지의 개별공시지가를 m\textsuperscript{2}당 1,126,850원으로 하향 조정해 줄 것을 내용으로 하는 이의신청을 하였다. 이에 대하여 甲은 이 사건 토지의 개별공시지가결정 시 표준지 선정에 문제가 있음을 발견하고, A시 부동산가격공시위원회의 심의를 거쳐 2018. 7. 1. 위 개별공시지가를 m\textsuperscript{2}당 2,380,000원으로 정정하여 결정ㆍ고시하였고, 동 결정서는 당일 乙에게 송달되었다. 丙은 2018. 6. 20. 위 이의신청과는 별개로 이 사건 토지의 개별공시지가를 m\textsuperscript{2}당 1,790,316원으로 수정해 달라는 취지의 행정심판을 청구하였고, B행정심판위원회는 2018. 8. 27. 이 사건 토지의 개별공시지가를 m\textsuperscript{2}당 2,000,000원으로 하는 변경재결을 하였고, 동 재결서 정본은 2018. 8. 30. 丙에게 송달되었다. 다음 물음에 답하시오. (30점) \par
\medskip\noindent 물음 1) 부동산 가격공시에 관한 법령상 개별공시지가의 정정사유에 관하여 설명하시오. (5점) \par
\medskip\noindent 물음 2) 위 사례에서 乙과 丙이 취소소송을 제기하려고 할 때, 소의 대상과 제소 기간의 기산일에 관하여 각각 설명하시오. (10점) \par
\medskip\noindent 물음 3) 한편, 丁은 A시의 개별공시지가 산정업무를 담당하고 있는 공무원이다. 丁은 개발예정지구인 C지역의 개별공시지가를 산정함에 있어 토지의 이용 상황을 잘못 파악하여 지가를 적정가격보다 훨씬 높은 가격으로 산정하였다. 이를 신뢰한 乙은 C지역의 담보가치가 충분하다고 믿고 그 토지에 근저당설정 등기를 마치고 수백억원의 투자를 하였지만, 결국 수십억원에 해당하는 큰 손해를 보았다. 이에 乙은 丁의 위법한 개별공시지가 산정으로 인하여 위 손해를 입었다고 주장하며, 국가배상소송을 제기하고자 한다. 동 소송에서 乙은 丁의 직무상 행위와 자신의 손해사이의 인과관계를 주장한다. 乙의 주장의 타당성에 관하여 개별공시지가제도의 입법목적을 중심으로 설명하시오. (15점) \par
\begin{framedtext}
참조 조문 \par
<부동산 가격공시에 관한 법률> \par
제12조(개별공시지가의 정정) 시장ㆍ군수 또는 구청장은 개별공시지가에 틀린 계산, 오기, 표준지 선정의 착오, 그 밖에 대통령령으로 정하는 명백한 오류가 있음을 발견한 때에는 지체 없이 이를 정정하여야 한다. \par
<행정소송법> \par
제19조(취소소송의 대상) 취소소송은 처분등을 대상으로 한다. 다만, 재결취소소송의 경우에는 재결 자체에 고유한 위법이 있음을 이유로 하는 경우에 한한다. \par
제20조(제소기간) ①취소소송은 처분등이 있음을 안 날부터 90일 이내에 제기하여야 한다. 다만, 제18조제1항 단서에 규정한 경우와 그 밖에 행정심판청구를 할 수 있는 경우 또는 행정청이 행정심판청구를 할 수 있다고 잘못 알린 경우에 행정심판청구가 있은 때의 기간은 재결서의 정본을 송달받은 날부터 기산한다. \par
②\textasciitilde{}③ <생략> \par
\end{framedtext}
\end{longquestion}
